\subsubsection{Intercambio masivo ordinario MASS-5}
\label{subsec:sd5_mass5}

Synaptix selecciona \texttt{MASS-5 1.0}, módulo propio permanente de la consola administrativa, para cargar y descargar información ordinaria, además de la migración inicial y de EXP-5. El cliente autorizado inicia, previsualiza, confirma, cancela y reanuda un trabajo sin ticket de soporte. CSV UTF-8 y JSON Lines son los formatos obligatorios; el paquete incorpora diccionario versionado, tipos, decimales, unidades, fechas UTC con zona de origen, nulos, IDs y relaciones. La descarga incluye datos y errores en esos formatos, no exige planilla propietaria ni se limita a la página visible. Los adjuntos conservan formato original, huella y referencia en manifiesto; los documentos no se interpretan como filas ejecutables.

\textbf{Origen, destino y autoridad.} El trabajo identifica usuario, finalidad, ámbito, dominio, esquema/versión, archivo/huella, sistema de origen, destino y operación solicitada. Descarga desde un corte del servicio propietario o su lectura versionada hacia la sesión autenticada del cliente; carga desde el archivo del cliente hacia las API de comando del dominio competente. MASS-5 no escribe directamente en tablas ni concede permisos nuevos: tanto la previsualización como cada aplicación y descarga verifican identidad, rol, unidad y autorización de campo. Deuda, documento tributario y pago externos se incorporan sólo como evidencia/solicitud a su adaptador y permanecen pendientes hasta conciliación con contabilidad; la carga no cambia su autoridad. Situación escolar, consentimiento, titularidad y acceso se someten a sus transiciones y custodios; importar una lista no confirma que una persona esté en el agua ni renueva un permiso. Se audita contenido antes/después mediante AUD-5 y se mantiene el UUID de correlación de INT-ACL.

\textbf{Validación previa y resultado por registro.} La recepción cifra el original y fija su huella. En staging sin efectos de negocio rechaza rutas de adjuntos fuera del paquete, archivos ejecutables y expansión superior al cupo; los parsers no ejecutan macros ni fórmulas. Después verifica cabecera/esquema, codificación, campos obligatorios, tipos, dominio de valores, unidades, fechas, unicidad, consentimiento y permisos. Un fallo de archivo impide confirmar y entrega diagnóstico de archivo; un error de fila produce número ordinal, ID lógico cuando exista, regla/código, campo, severidad y explicación corregible, sin copiar secretos ni datos ajenos al diagnóstico. Verifica referencias contra el corte y contra otros registros del lote; ordena las dependencias en grupos atómicos mínimos. Padre e hijos que deban existir juntos se aplican en una transacción del servicio propietario o se rechazan como grupo. Referencias desconocidas, ciclos no admisibles y autorización insuficiente impiden ese grupo, conservando los independientes como aplicables. La previsualización presenta cantidades válidas, inválidas, dependientes y advertidas; confirmar procesamiento parcial exige elección expresa del cliente y sólo aplica grupos válidos. No compensa una fila inválida con un total monetario aparentemente correcto. La integridad se vuelve a validar al aplicar: una referencia borrada, versión cambiada o permiso revocado desde la previsualización produce conflicto por fila, sin sobrescribir el cambio concurrente.

\textbf{Reanudación e idempotencia.} Cada fila/grupo conserva clave formada por cliente, origen, trabajo, ordinal/ID y operación, huella del contenido, versión esperada y estado validado, pendiente, aplicado, rechazado, en conflicto o resultado desconocido. El servicio propietario persiste esa clave y resultado en la misma transacción del efecto y outbox. Reenviar el mismo contenido devuelve el resultado anterior; reutilizar clave con otra huella se rechaza. Un timeout sin respuesta no se cuenta como fallo definitivo ni se aplica de nuevo a ciegas: consulta el resultado durable por clave. Una corrección se entrega como nueva revisión vinculada al rechazo; conserva los aplicados y no los duplica. El resumen cumple total recibido = aplicado + rechazado + conflicto + pendiente/resultado desconocido, sin contar filas de cabecera; adjunta informe por registro y por grupo. Un trabajo cancelado deja constancia de lo ya aplicado; no promete deshacer efectos contables o externos. Descarga masiva fija corte/filtros/orden y manifiesto global, divide en partes reanudables y verifica cantidades/IDs/huellas; cambio de permisos detiene las partes aún no entregadas.

\textbf{Capacidad de carga y aceptación contractual.}
Se mantiene una carga ordinaria activa y una pendiente por ámbito, partes de hasta 10.000 filas o 20 MiB y hasta 100 partes por trabajo; un conjunto mayor se encadena por su mismo manifiesto, sin truncarlo. Los grupos atómicos conservan integridad de padres e hijos y consumen recursos por todos sus registros, sin fraccionar una transacción para aparentar una tasa mayor.
La aceptación exige al menos 10.000 registros por minuto con registros representativos y validación activa, equivalentes a 166,67 registros por segundo. Se individualiza para cada ensayo el alcance de validación, reglas, persistencia, auditoría y eventos, errores por fila y relaciones. El perfil preliminar de diez validaciones y dos aplicaciones por segundo, y su crecimiento treinta/seis, no acredita esta exigencia y se retira como perfil de recepción contractual. La estimación de 5.000 segundos para aplicar 10.000 filas pertenecía a ese perfil preliminar y no es la duración ofrecida del perfil contractual.
La consola conserva progreso, estimación, edad, dueño, motivo de pausa y plazo reestimado. Operación crítica, descarga, informes y adjuntos mantienen reservas explícitas; no se acepta un cargador que nunca termina bajo la mezcla contractual. Se ensayan padres/hijos, filas válidas e inválidas, permiso revocado, concurrencia, timeout tras commit, cancelación y reinicio, con resultado por fila, aplicación parcial autorizada y reejecución sin duplicados.

\subsubsection{Calendario y entrega automática de todos los informes ENV-5}
\label{subsec:sd5_env5}

Se selecciona \texttt{ENV-5 1.0}, módulo propio de calendario y envío integrado en la consola, para \emph{todo informe} producido por la solución: plantillas oficiales de los doce dominios, listados e informes operacionales, financieros, ambientales, escolares, de auditoría/soporte y cualquier diseño creado por el cliente en ANAL-5. Los 44 indicadores C07 no delimitan este universo. REG-INF-5 es el registro versionado obligatorio de informe, propietario, esquema de salida, parámetros, corte, sensibilidad, exportador y permisos. No se publica un informe nuevo ni una revisión sin inscribirlo y comprobar su exportación y calendario. Representaciones gráficas se acompañan de series/datos y especificación del gráfico en JSON; texto, anexos e imágenes incluyen originales y manifiesto de relaciones, sin hacer de PDF la única exportación. CSV UTF-8 y JSON documentados cubren todo informe; Parquet es complemento opcional. Un extracto de contabilidad mantiene su autoridad y frescura declarada; programarlo no cierra RT-05.29.

\textbf{Calendario.} Un Autor autorizado configura destinatarios nominales, periodicidad diaria/semanal/mensual, fecha/hora local, zona \texttt{America/Santiago}, vigencia, filtros/corte y formato. Permite una fecha única y varias horas de un día. Se almacena UTC más regla local: una hora repetida se ejecuta una vez con el primer offset y una inexistente pasa al primer instante local válido, dejando evidencia; el fin de mes usa el último día existente. El Autor ve próximos disparos y confirma esa regla. Una ejecución tiene ID único programa--versión--instante previsto; el planificador durable evita duplicarla al reiniciar y recupera los disparos vencidos, etiquetando fecha prevista, real y retraso, sin reemplazarlos por el próximo período. Cambios/cancelación no alteran informes ya emitidos. Antes de extraer y antes de entregar se revalidan autor, destinatario, ámbito y finalidad. El contenido se genera por universo autorizado de cada destinatario, no se envía un archivo amplio confiando en que éste no lo abra. Revocación suspende ese destinatario, registra motivo y alerta al custodio; modificar una lista o compartir un diseño no otorga autorización.

\textbf{Entrega del archivo por canal propio.} El canal seleccionado es la bandeja documental nominativa del portal web, accesible a destinatarios internos y externos autorizados con identidad individual y autenticación reforzada. ENV-5 transmite automáticamente las partes cifradas del informe y su manifiesto desde su outbox a la bandeja de cada destinatario; allí persiste una copia del contenido dirigida a esa identidad. No se registra entrega por guardar el original de generación, publicar un enlace común, mostrar disponibilidad al autor o enviar un aviso por correo. El receptor propio verifica huellas, número de partes y ámbito, confirma persistencia durable de la copia y devuelve recibo autenticado con ejecución, destinatario, manifiesto e instante; sólo ese recibo marca entregado. La apertura/descarga por el humano es un estado posterior separado, no requisito para el envío automático. La copia se cifra por entrega, con envoltura asociada al receptor; el servicio de lectura exige permiso actual, no una URL portadora, y audita descarga. Los avisos de consola pueden acompañar esa entrega, sin sustituirla. No depende de SMS, WhatsApp, SES ni del servicio de emergencia. El cliente administra programas y destinatarios sin intervención de Synaptix. No se habilitan destinatarios anónimos ni cuentas compartidas para escolares, deuda o datos sensibles.

\textbf{Fallos y retención.} Generación, transmisión, recibo y apertura tienen estados separados. Una pérdida de recibo obliga a consultar la entrega por clave ejecución--destinatario--huella y a reanudar sólo partes faltantes; distinta huella con misma clave es conflicto. Fallos transitorios reintentan a 1, 5, 15 y 60 minutos, después cada hora hasta siete días, con aviso al custodio desde el primer fallo y escalamiento a Operación al superar quince minutos. El envío no desaparece al agotarse reintentos: permanece fallido con causa/recibo faltante y exige corrección y reenvío auditado. La caída de WAN mantiene los disparos en cola; retraso se registra como tal, sin declarar entrega puntual. Se conservan informes y bandejas por siete días desde su entrega, visibles con fecha de vencimiento; el programa y sus recibos sin payload siguen la retención de auditoría. Los originales de negocio mantienen su plazo propio. El acceso revocado bloquea lectura y envíos posteriores; la revocación no modifica copias ya descargadas por el cliente.

\label{subsec:sd5_env5_recursos}
\paragraph{Recursos comunes y revisión del perfil MASS-5/ENV-5:}
El planificador y los comandos HTTP se cuentan una sola vez en Marina-Admin. Los trabajos de MASS-5 y ENV-5 requieren recursos de ejecución explícitos y reserva conciliada con analítica, exportación, reconexión y negocio. Con las hipótesis de 0,020 segundos CPU por validación, 0,010 por aplicación, 0,020 de serialización y 0,040 de transferencia, una referencia de 167 registros por segundo exige
\[
C_{\mathrm{lote}}=167(0{,}020+0{,}010)+0{,}020+0{,}040=5{,}07\ \mathrm{vCPU}.
\]
Es un piso aritmético previo a margen, agentes, tráfico y falla, no rendimiento demostrado. Una tarea de 1 vCPU como la descrita por SD4 no ejecuta por sí sola ese trabajo al ritmo exigido bajo estas hipótesis. Las seis/doce tareas de los tres núcleos tampoco constituyen una reserva dedicada a MASS/ENV. Antes de cerrar el perfil se seleccionan ejecutor, concurrencia, CPU, memoria, I/O, caudales y recuperación en Santiago, conciliados con SD4/T11. Las cuotas son globales y no se multiplican por controlador, ámbito o destinatario. La toma por un sustituto mantiene claves idempotentes y resultados durables.

Se propone un buffer compartido de hasta 128 MiB y temporal de 256 MiB por controlador para ANAL-5, MASS-5, ENV-5 y EXP-5, sin sumar cuatro reservas independientes. El heap se dimensiona dentro de la memoria total de la tarea, descontando las otras necesidades del proceso. Las partes se transfieren a objetos cifrados y su staging se borra al comprobar persistencia. Se reserva capacidad \emph{adicional} de 170 GiB de objetos iniciales, 510 GiB con crecimiento, distinta de los subtotales de datos/journal: 10 GiB para cargas retenidas hasta treinta días para corrección y 160 GiB para originales de informes y copias nominativas de siete días (30 y 480 GiB con crecimiento). Como mezcla inicial se consideran cien disparos diarios de 20 MiB retenidos siete días y hasta diez destinatarios por disparo, además del original: $100\times7\times20\times11/1024=150,39$ GiB. Los 9,61 GiB restantes de la reserva de informes son margen para estados transitorios, no una autorización de acumulación ilimitada. Los originales y partes de envíos aún pendientes o fallidos conservan custodia hasta resolverlos; antes de superar la envolvente se amplía almacenamiento, sin purgar esos archivos para fingir espacio disponible ni marcar entrega sin recibo. Mayor cantidad de destinatarios, tamaños o frecuencia se divide en partes pero no reduce el volumen: se amplía esa reserva antes de admitir el calendario, y se valoriza en OE bajo los recursos existentes. A 2 MiB/s, transmitir un informe de 20 MiB a diez bandejas necesita al menos cien segundos, además de extracción y recibos. MASS-5 y ENV-5 comparten coordinación durable y cuotas globales de validación, aplicación y caudal. La cantidad de workers pesados y su distribución se fija con el perfil contractual y sus recursos de SD4/T11; la referencia anterior de uno/tres workers no es una configuración definitiva. Las cuotas no se multiplican por las tareas, ámbitos ni destinatarios. El receptor persiste por streaming fuera del heap y vuelve a verificar el manifiesto antes del recibo. ENV-5 admite inicialmente hasta 100 disparos diarios y diez destinatarios por disparo como mezcla de capacidad, no como límite contractual de la función: cantidades superiores requieren ampliar recurso antes de activar el programa y quedan visibles, nunca se silencian ni se anuncian como enviadas. Los recibos, programas y errores forman parte de AUD-5; su ocupación debe entrar en el dimensionamiento integral todavía pendiente de RT-05.05/13/14, sin presentar los 170 GiB como tamaño total de la solución.

La aceptación recorrerá el registro completo REG-INF-5, incluido un informe nuevo del cliente y uno de cada dominio, exportará formatos abiertos y leerá contenido/diccionario en herramienta independiente. Programará calendario, medianoche, horario repetido/inexistente y fin de mes; verificará para cada destinatario recibido contra huella/partes y versión/filtros, y diferenciará entregado de abierto. Incluirá revocación, reinicio, corte WAN, pérdida de recibo, reintento, destinatario ajeno y cola concurrente con peak, crecimiento y zona retirada. La recepción exige completar los envíos sin perjudicar las funciones críticas; se amplía capacidad si no cumple. No se afirma que estas pruebas ni entregas se hayan ejecutado (BT RT-05.28).

\subsubsection{Presupuesto de actualidad y brechas de RT-05.29}
\label{subsec:sd5_actualidad29}

La medición parte de la ocurrencia física, no de la recepción en nube: $L=D_{captura}+D_{enlace}+D_{cola}+D_{aplicacion}+D_{vista}$, más incertidumbre de reloj registrada. Ocupación mantiene un máximo de sesenta segundos desde la ocurrencia física hasta la primera vista autorizada actualizada. El presupuesto de diseño asigna cinco segundos al registro, quince al enlace, diez a la cola, quince a persistencia/proyección, tres a vista y hasta cinco a incertidumbre del reloj: cincuenta y tres segundos, con siete de margen. Se comprueba además el límite propio de tres segundos desde aceptación por la autoridad, conforme a 5.1.2.2 y Tabla~5.22. El ensayo verifica ambos; no basta medir desde el acuse si el evento esperó desconectado. Son presupuestos de diseño, no tiempos medidos. Para consumo, adquisición con espera máxima de un minuto y publicación con espera máxima de quince minutos dejan hasta 44 minutos para transmisión, validación, consolidación, vista e incertidumbre de reloj, dentro de una hora desde la ocurrencia: uno más quince más cuarenta y cuatro son sesenta minutos; una lectura inexistente no se sustituye por cero. Ambiente fija a las 02:00 locales el cierre diario del día anterior después de conciliar fuentes y declara los faltantes del día; no usa un informe parcial para certificar consolidación completa.

En cuenta, el contrato necesario exige lectura de la autoridad contable en la consulta y publicación por cada cambio de cargo/pago/anulación, con ID/version y secuencia, sin batch como vía ordinaria. INT-ACL y MASS-5 conservan eventos/evidencia, pero todavía no existe en las fuentes un contrato de una interfaz contable efectiva que demuestre esos mecanismos. Antes de habilitar se debe seleccionar y documentar endpoint/operaciones, identidad admitida, lectura coherente, eventos o alternativa sin consolidación periódica, continuidad y reconciliación con su propietario; un emulador no acredita esa disponibilidad. El extracto periódico se conserva sólo como contingencia identificada.

\paragraph{Recorrido del estado de cuenta.}
La contabilidad externa conserva la autoridad sobre documento tributario, pago, anulación y saldo. FAC consulta el estado vigente a la autoridad contable y recibe desde ella cada cambio confirmado de cargo, pago o anulación. El adaptador conserva la dirección de ambos intercambios: solicitud FAC--contabilidad y respuesta o cambio contabilidad--FAC. El contrato lógico exige estos recorridos sin utilizar un extracto periódico como ruta ordinaria. Es una especificación de la solución; no identifica una API de tercero ya habilitada.

Una respuesta o cambio identifica cuenta, armador o contrato, documento de origen, operación, versión o secuencia de fuente, instante del hecho, instante de confirmación contable, moneda, importe, estado y correlación. INT-ACL conserva entrada, respuesta original y resultado durable. FAC deduplica por origen e identificador, valida referencias y publica su actualización después de confirmar la transacción. La consulta autorizada conserva el corte de la fuente y verifica su vigencia; una proyección cuya versión no se pueda confirmar se muestra como saldo no comprobado, no como cero ni como saldo actual. Diferencias y respuestas desconocidas quedan identificadas para reconciliación sin repetir un cargo o pago.

La aceptación recorre cargo, pago, anulación, respuesta perdida, duplicado, cambio fuera de orden y recuperación. Mide desde la confirmación del hecho en contabilidad hasta su primera visualización autorizada. La ficha efectiva de la contraparte debe demostrar consulta y cambios con esta semántica; un emulador o extracto periódico permite ensayar el adaptador, pero no acredita el requisito de tiempo real.

\paragraph{Recorrido de la situación escolar.}
ESC mantiene una autoridad única en el recinto para participación, embarcación, monitor, regreso y entrega. Portería, escuela y monitor utilizan identificadores de alumno, sesión y participación comunes. Cada salida, cambio, regreso o retiro incluye operación, actor e informante, dispositivo, versión esperada, instante de ocurrencia, recepción y evidencia. El servicio comprueba permiso vigente y transición; persiste participación, auditoría y outbox en la misma transacción y emite un acuse con la versión aceptada. La difusión notifica el cambio y cada consulta recupera esa versión del servicio propietario, sin esperar el lote analítico diario.

Un cambio de nave o monitor cierra el intervalo anterior y abre el siguiente; un retiro sólo se confirma con identidad y facultad comprobadas y entrega efectiva. El cierre de clase requiere regreso individual de quienes efectivamente entraron al agua. La caída de WAN con LAN operativa conserva el recorrido contra la autoridad local. La captura en un equipo sin camino hacia esa autoridad queda identificada como no compartida; no se presenta como situación escolar confirmada.

La prueba registra, para portería, escuela y monitor, la versión del evento y el instante de primera vista autorizada, junto con incertidumbre de reloj. Incluye salida, cambio, regreso, retiro, duplicado, permiso vencido y partición. Sólo se declara conforme RT-05.29 cuando el recorrido mantiene el estado compartido exigido; la captura aislada y el aviso radial no sustituyen esa comprobación.

Mostrar lista antigua, indeterminado, radio o ingresar posteriormente no satisface el tiempo real. Calidad conservará instantes de ocurrencia, recepción, commit, publicación y primera vista autorizada con incertidumbre de reloj para cada fuente, y reportará máximo y cada incumplimiento; no sólo promedio/P95. Ensayará el viernes crítico, eventos de fuentes simultáneamente aisladas, saldo externo y autonomías de 24/12/8 h cruzadas con antigüedad del permiso. RT-05.29 permanece con desarrollo por completar mientras estos mecanismos no estén resueltos, sin rebajar los cinco complementos C07 a las cuatro horas generales de BT.
